\addChapter{\textit{ABSTRACT}}
\vspace{0.1cm}

{
    \singlespacing
    \begin{center}
        \textbf{\textit{ABSTRACT}}
    \end{center}
    \vspace{0.5cm}

    {\frenchspacing
        \textit{
            The process of shooting film scenes is often not carried out sequentially according to the storyline.
            One way to determine the shooting order of scenes depends on logistical efficiency, such as transportation costs and talent fees.
            An unoptimized shooting sequence can lead to increased production costs.
            This study addresses the problem of determining the optimal shooting order of film scenes as a form of the Traveling Salesman Problem (TSP), solved using the Particle Swarm Optimization (PSO) algorithm.
            A total of 60 scenes are represented in the form of a complete graph, where each scene is modeled as a vertex and the shooting cost between scenes as weighted edges.
            The shooting sequence starts and ends at scene 4 in accordance with the TSP concept.
            The data were processed through normalization, fitness calculation, and PSO iteration stages to obtain a scene sequence that produces the minimum total cost.
            The results show that the application of the PSO algorithm can provide a shooting sequence solution with a minimum total cost of Rp$7,794,359.48$.
            The optimized sequence obtained is: 4-1C-22B-27-10-12B-32A-23B-22C-2A-34B-23A-33B-33A-19-31-1E-7A-32B-24-16B-22A-3-35B-14B-9A-28-12A-29A-21B-35A-
            37-6A-25-8-7B-30A-26-16A-11-2B-19B-1D-13-5-34-6B-20-9B\newline
            -15-17-16D-18-29B-30B-16C-30C-14A-21A-36-4.
            Thus, the PSO algorithm provides as an alternative solution for determining the shooting order of scenes in film production.
        }

        \vspace*{0.5cm}
        \noindent \textbf{\textit{Keywords:}} \textit{graph}, \textit{particle swarm optimization} (PSO), \textit{shooting order of scenes}, \textit{traveling salesman problem} (TSP)

    }
}

\pagebreak